\IfFileExists{stacks-project.cls}{%
\documentclass{stacks-project}
}{%
\documentclass{amsart}
}
\usepackage{amssymb}

% For dealing with references we use the comment environment
\usepackage{verbatim}
\usepackage{fancyvrb}
\usepackage{pmboxdraw}
\newenvironment{reference}{\comment}{\endcomment}
%\newenvironment{reference}{}{}
\newenvironment{slogan}{\comment}{\endcomment}
\newenvironment{history}{\comment}{\endcomment}

% For source figures retained as assets
\usepackage{graphicx}

% For commutative diagrams we use Xy-pic
\usepackage[all]{xy}

% We use 2cell for 2-commutative diagrams.
\xyoption{2cell}
\UseAllTwocells

% We use multicol for the list of chapters between chapters
\usepackage{multicol}

% This is generally recommended for better output
\usepackage{lmodern}
\usepackage[T1]{fontenc}
\usepackage{textalpha}

% For cross-file-references
% We prefer the snapshot tmp/xref/#2.aux,
% which latexmkrc replaces atomically, so
% that documents compiled simultaneously
% never read a half-written .aux file.
% The last argument keeps links to #2.pdf.
\usepackage{xr-hyper}
\newcommand{\maybeexternaldocument}[2]{%
  \IfFileExists{#2.aux}{%
    \externaldocument[#1]{#2}%
  }{%
    \PackageWarning{stack}{Missing #2.aux; external references with prefix #1 unavailable}%
  }%
}
\let\externaldocumentorig\externaldocument
\renewcommand{\externaldocument}[2][]{%
  \IfFileExists{tmp/xref/#2.aux}{%
    \externaldocumentorig[#1]{tmp/xref/#2}[#2.pdf]%
  }{%
  \IfFileExists{#2.aux}{%
    \externaldocumentorig[#1]{#2}%
  }{%
    \PackageWarning{stack}{Missing #2.aux; external references with prefix #1 unavailable}%
  }%
  }%
}

% Package for hypertext links:
\usepackage{hyperref}

% For any local file, say "hello.tex" you want to link to please
% use \externaldocument[hello-]{hello}
\externaldocument[introduction-]{introduction}
\externaldocument[conventions-]{conventions}
\externaldocument[sets-]{sets}
\externaldocument[categories-]{categories}
\externaldocument[topology-]{topology}
\externaldocument[sheaves-]{sheaves}
\externaldocument[sites-]{sites}
\externaldocument[stacks-]{stacks}
\externaldocument[fields-]{fields}
\externaldocument[algebra-]{algebra}
\externaldocument[brauer-]{brauer}
\externaldocument[homology-]{homology}
\externaldocument[derived-]{derived}
\externaldocument[simplicial-]{simplicial}
\externaldocument[more-algebra-]{more-algebra}
\externaldocument[smoothing-]{smoothing}
\externaldocument[modules-]{modules}
\externaldocument[sites-modules-]{sites-modules}
\externaldocument[injectives-]{injectives}
\externaldocument[cohomology-]{cohomology}
\externaldocument[sites-cohomology-]{sites-cohomology}
\externaldocument[dga-]{dga}
\externaldocument[dpa-]{dpa}
\externaldocument[sdga-]{sdga}
\externaldocument[hypercovering-]{hypercovering}
\externaldocument[schemes-]{schemes}
\externaldocument[constructions-]{constructions}
\externaldocument[properties-]{properties}
\externaldocument[morphisms-]{morphisms}
\externaldocument[coherent-]{coherent}
\externaldocument[divisors-]{divisors}
\externaldocument[limits-]{limits}
\externaldocument[varieties-]{varieties}
\externaldocument[topologies-]{topologies}
\externaldocument[descent-]{descent}
\externaldocument[perfect-]{perfect}
\externaldocument[more-morphisms-]{more-morphisms}
\externaldocument[flat-]{flat}
\externaldocument[groupoids-]{groupoids}
\externaldocument[more-groupoids-]{more-groupoids}
\externaldocument[etale-]{etale}
\externaldocument[chow-]{chow}
\externaldocument[intersection-]{intersection}
\externaldocument[pic-]{pic}
\externaldocument[weil-]{weil}
\externaldocument[adequate-]{adequate}
\externaldocument[dualizing-]{dualizing}
\externaldocument[duality-]{duality}
\externaldocument[discriminant-]{discriminant}
\externaldocument[derham-]{derham}
\externaldocument[local-cohomology-]{local-cohomology}
\externaldocument[algebraization-]{algebraization}
\externaldocument[curves-]{curves}
\externaldocument[resolve-]{resolve}
\externaldocument[models-]{models}
\externaldocument[functors-]{functors}
\externaldocument[equiv-]{equiv}
\externaldocument[pione-]{pione}
\externaldocument[etale-cohomology-]{etale-cohomology}
\externaldocument[proetale-]{proetale}
\externaldocument[relative-cycles-]{relative-cycles}
\externaldocument[more-etale-]{more-etale}
\externaldocument[trace-]{trace}
\externaldocument[crystalline-]{crystalline}
\externaldocument[spaces-]{spaces}
\externaldocument[spaces-properties-]{spaces-properties}
\externaldocument[spaces-morphisms-]{spaces-morphisms}
\externaldocument[decent-spaces-]{decent-spaces}
\externaldocument[spaces-cohomology-]{spaces-cohomology}
\externaldocument[spaces-limits-]{spaces-limits}
\externaldocument[spaces-divisors-]{spaces-divisors}
\externaldocument[spaces-over-fields-]{spaces-over-fields}
\externaldocument[spaces-topologies-]{spaces-topologies}
\externaldocument[spaces-descent-]{spaces-descent}
\externaldocument[spaces-perfect-]{spaces-perfect}
\externaldocument[spaces-more-morphisms-]{spaces-more-morphisms}
\externaldocument[spaces-flat-]{spaces-flat}
\externaldocument[spaces-groupoids-]{spaces-groupoids}
\externaldocument[spaces-more-groupoids-]{spaces-more-groupoids}
\externaldocument[bootstrap-]{bootstrap}
\externaldocument[spaces-pushouts-]{spaces-pushouts}
\externaldocument[spaces-chow-]{spaces-chow}
\externaldocument[groupoids-quotients-]{groupoids-quotients}
\externaldocument[spaces-more-cohomology-]{spaces-more-cohomology}
\externaldocument[spaces-simplicial-]{spaces-simplicial}
\externaldocument[spaces-duality-]{spaces-duality}
\externaldocument[formal-spaces-]{formal-spaces}
\externaldocument[restricted-]{restricted}
\externaldocument[spaces-resolve-]{spaces-resolve}
\externaldocument[formal-defos-]{formal-defos}
\externaldocument[defos-]{defos}
\externaldocument[cotangent-]{cotangent}
\externaldocument[examples-defos-]{examples-defos}
\externaldocument[algebraic-]{algebraic}
\externaldocument[examples-stacks-]{examples-stacks}
\externaldocument[stacks-sheaves-]{stacks-sheaves}
\externaldocument[criteria-]{criteria}
\externaldocument[artin-]{artin}
\externaldocument[quot-]{quot}
\externaldocument[stacks-properties-]{stacks-properties}
\externaldocument[stacks-morphisms-]{stacks-morphisms}
\externaldocument[stacks-limits-]{stacks-limits}
\externaldocument[stacks-cohomology-]{stacks-cohomology}
\externaldocument[stacks-perfect-]{stacks-perfect}
\externaldocument[stacks-introduction-]{stacks-introduction}
\externaldocument[stacks-more-morphisms-]{stacks-more-morphisms}
\externaldocument[stacks-geometry-]{stacks-geometry}
\externaldocument[moduli-]{moduli}
\externaldocument[moduli-curves-]{moduli-curves}
\externaldocument[examples-]{examples}
\externaldocument[exercises-]{exercises}
\externaldocument[guide-]{guide}
\externaldocument[desirables-]{desirables}
\externaldocument[coding-]{coding}
\externaldocument[obsolete-]{obsolete}
\externaldocument[fdl-]{fdl}
\externaldocument[index-]{index}

% Heap Project documents
\externaldocument[k3-]{k3}
\externaldocument[torelli-]{torelli}
\externaldocument[linear-algebra-]{linear-algebra}
\externaldocument[groups-]{groups}
\externaldocument[complex-analysis-]{complex-analysis}
\externaldocument[functional-analysis-]{functional-analysis}
\externaldocument[categories-]{categories}
\externaldocument[commutative-algebra-]{commutative-algebra}
\externaldocument[ringed-spaces-]{ringed-spaces}
\externaldocument[homological-algebra-]{homological-algebra}
\externaldocument[lie-algebras-]{lie-algebras}
\externaldocument[representation-theory-]{representation-theory}
\externaldocument[smooth-manifolds-]{smooth-manifolds}
\externaldocument[differential-topology-]{differential-topology}
\externaldocument[algebraic-topology-]{algebraic-topology}
\externaldocument[differential-geometry-]{differential-geometry}
\externaldocument[algebraic-geometry-]{algebraic-geometry}
\externaldocument[symplectic-geometry-]{symplectic-geometry}
\externaldocument[complex-geometry-]{complex-geometry}
\externaldocument[hodge-theory-]{hodge-theory}
\externaldocument[%
variations-of-hodge-structures-]{%
variations-of-hodge-structures}
\externaldocument[%
homological-mirror-symmetry-]{%
homological-mirror-symmetry}
\maybeexternaldocument{deformations-}{deformations}
\externaldocument[vertex-algebras-]{vertex-algebras}
\externaldocument[springer-]{springer}
\externaldocument[infty-categories-]{infty-categories}
\externaldocument[seminars-differential-geometry-]{%
seminars-differential-geometry}
\externaldocument[seminars-complex-geometry-]{%
seminars-complex-geometry}
\externaldocument[seminars-algebraic-geometry-]{%
seminars-algebraic-geometry}
\externaldocument[seminars-analysis-pde-]{%
seminars-analysis-pde}
\externaldocument[surfaces-]{surfaces}
\externaldocument[birational-maps-commutative-algebra-]{birational-maps-commutative-algebra}
\externaldocument[cremona-transformations-]{cremona-transformations}
\externaldocument[derived-categories-of-sheaves-]{derived-categories-of-sheaves}
\externaldocument[geometric-invariant-theory-]{geometric-invariant-theory}
\externaldocument[geometric-stability-conditions-and-group-actions-]{geometric-stability-conditions-and-group-actions}
\externaldocument[moduli-spaces-of-sheaves-]{moduli-spaces-of-sheaves}
\externaldocument[bridgeland-stability-general-theory-]{bridgeland-stability-general-theory}
\externaldocument[hilbert-schemes-]{hilbert-schemes}
\externaldocument[noncommutative-abelian-surfaces-and-kummer-type-hyperkahler-varieties-]{noncommutative-abelian-surfaces-and-kummer-type-hyperkahler-varieties}
\externaldocument[mumford-tate-groups-in-hodge-theory-]{mumford-tate-groups-in-hodge-theory}
\externaldocument[hodge-birational-atoms-2026-]{hodge-birational-atoms-2026}
\externaldocument[xxii-egd-2026-]{%
xxii-egd-2026}
\externaldocument[pde-]{pde}
\externaldocument[probability-]{probability}

\externaldocument[linguistics-]{linguistics}

\externaldocument[%
16th-alga-meeting-2026-]{%
16th-alga-meeting-2026}

\externaldocument[%
iv-jornadas-lefschetz-2026-]{%
iv-jornadas-lefschetz-2026}

\externaldocument[reading-the-masters-]{%
reading-the-masters}

\externaldocument[real-analysis-]{%
real-analysis}

\externaldocument[optimal-transport-]{%
optimal-transport}

\externaldocument[history-]{history}

\externaldocument[seiberg-witten-]{seiberg-witten}

\externaldocument[history-philosophy-biology-]{history-philosophy-biology}

\externaldocument[classical-mechanics-]{classical-mechanics}
\externaldocument[electromagnetism-]{electromagnetism}
\externaldocument[waves-and-optics-]{waves-and-optics}
\externaldocument[wif-]{waves-in-fluids-iv}
\externaldocument[thermal-physics-]{thermal-physics}
\externaldocument[quantum-mechanics-]{quantum-mechanics}
\externaldocument[special-relativity-]{special-relativity}
\externaldocument[physics-]{physics}

\externaldocument[chemistry-]{chemistry}

\externaldocument[molecular-cell-biology-]{molecular-cell-biology}

\externaldocument[microbiology-]{microbiology}

\externaldocument[protistology-]{protistology}

\externaldocument[mycology-]{mycology}

\externaldocument[botany-]{botany}

\externaldocument[zoology-]{zoology}

\externaldocument[animal-histology-]{animal-histology}

\externaldocument[animal-physiology-]{animal-physiology}

\externaldocument[systematics-biogeography-]{systematics-biogeography}

\externaldocument[ecology-natural-resources-]{ecology-natural-resources}

\externaldocument[earth-sciences-]{earth-sciences}

\externaldocument[biostatistics-]{biostatistics}

\externaldocument[mathematical-systems-biology-]{mathematical-systems-biology}

\externaldocument[cimpa-el-salvador-2024-]{cimpa-el-salvador-2024}

\externaldocument[anthropology-]{anthropology}

\externaldocument[yang-mills-theory-]{yang-mills-theory}

\externaldocument[seminars-algebra-]{seminars-algebra}

\externaldocument[seminars-physics-]{seminars-physics}

\externaldocument[seminars-computing-]{seminars-computing}

\externaldocument[oktobermat-2026-]{oktobermat-2026}

\externaldocument[saet-2026-]{saet-2026}

\externaldocument[seminars-symplectic-geometry-]{seminars-symplectic-geometry}

\externaldocument[seminars-dynamics-]{seminars-dynamics}

\externaldocument[oktobermat-2024-]{oktobermat-2024}

\externaldocument[quantum-groups-]{quantum-groups}

\externaldocument[legal-2027-]{legal-2027}

\externaldocument[derived-categories-]{derived-categories}

% Theorem environments.
%
\theoremstyle{plain}
\newtheorem{theorem}[subsection]{Theorem}
\newtheorem{proposition}[subsection]{Proposition}
\newtheorem{lemma}[subsection]{Lemma}
\newtheorem{corollary}[subsection]{Corollary}

\theoremstyle{definition}
\newtheorem{definition}[subsection]{Definition}
\newtheorem{example}[subsection]{Example}
\newtheorem{exercise}[subsection]{Exercise}
\newtheorem{situation}[subsection]{Situation}

\theoremstyle{remark}
\newtheorem{remark}[subsection]{Remark}
\newtheorem{remarks}[subsection]{Remarks}

\numberwithin{equation}{subsection}

% Macros
%
\def\lim{\mathop{\mathrm{lim}}\nolimits}
\def\colim{\mathop{\mathrm{colim}}\nolimits}
\def\Spec{\mathop{\mathrm{Spec}}}
\def\Hom{\mathop{\mathrm{Hom}}\nolimits}
\def\Ext{\mathop{\mathrm{Ext}}\nolimits}
\def\SheafHom{\mathop{\mathcal{H}\!\mathit{om}}\nolimits}
\def\SheafExt{\mathop{\mathcal{E}\!\mathit{xt}}\nolimits}
\def\Sch{\mathit{Sch}}
\def\Mor{\mathop{\mathrm{Mor}}\nolimits}
\def\Ob{\mathop{\mathrm{Ob}}\nolimits}
\def\Sh{\mathop{\mathit{Sh}}\nolimits}
\def\NL{\mathop{N\!L}\nolimits}
\def\CH{\mathop{\mathrm{CH}}\nolimits}
\def\proetale{{pro\text{-}\acute{e}tale}}
\def\etale{{\acute{e}tale}}
\def\QCoh{\mathit{QCoh}}
\def\Ker{\mathop{\mathrm{Ker}}}
\def\Im{\mathop{\mathrm{Im}}}
\def\Coker{\mathop{\mathrm{Coker}}}
\def\Coim{\mathop{\mathrm{Coim}}}

% Boxtimes
%
\DeclareMathSymbol{\boxtimes}{\mathbin}{AMSa}{"02}

%
% Macros for moduli stacks/spaces
%
\def\QCohstack{\mathcal{QC}\!\mathit{oh}}
\def\Cohstack{\mathcal{C}\!\mathit{oh}}
\def\Spacesstack{\mathcal{S}\!\mathit{paces}}
\def\Quotfunctor{\mathrm{Quot}}
\def\Hilbfunctor{\mathrm{Hilb}}
\def\Curvesstack{\mathcal{C}\!\mathit{urves}}
\def\Polarizedstack{\mathcal{P}\!\mathit{olarized}}
\def\Complexesstack{\mathcal{C}\!\mathit{omplexes}}
% \Pic is the operator that assigns to X its picard group, usage \Pic(X)
% \Picardstack_{X/B} denotes the Picard stack of X over B
% \Picardfunctor_{X/B} denotes the Picard functor of X over B
\def\Pic{\mathop{\mathrm{Pic}}\nolimits}
\def\Picardstack{\mathcal{P}\!\mathit{ic}}
\def\Picardfunctor{\mathrm{Pic}}
\def\Deformationcategory{\mathcal{D}\!\mathit{ef}}


\begin{document}

\title{Triangulated spheres and polytopality}
\maketitle

\phantomsection
\label{section-phantom}

\tableofcontents

\section{Convex polytopes and their
boundaries}
\label{section-convex-polytopes}

\noindent
This chapter collects notes on convex
polytopes and on triangulated spheres and
manifolds. The basic definitions are
standard; see \cite{ziegler-polytopes}.
Polytopes in the skeletal sense, where a
polygon is a cycle of edges with no interior,
are the subject of the chapter Skeletal
polytopes, Section
\ref{skeletal-polytopes-section-skeletal}.

A convex polytope is the convex hull of a
finite set of points in $\mathbb{R}^n$. Its
dimension is the dimension of the affine
space it spans. A face of a convex polytope
$P$ is the set of points of $P$ where a
linear function attains its maximum on $P$.
The faces of dimension $0$ are the vertices,
and the faces of dimension one less than $P$
are the facets.

A convex polytope is simplicial when every
facet is a simplex. In that case every proper
face is a simplex, and the proper faces form
a simplicial complex on the set of vertices:
a set of vertices is a face of the complex
when it is the set of vertices of a proper
face of $P$. This is the boundary complex of
$P$. If $P$ has dimension $d+1$, its boundary
is homeomorphic to the sphere of dimension
$d$, so the boundary complex is a
triangulation of that sphere.

Simplicial complexes are defined in
Stanley-Reisner schemes, Section
\ref{stanley-reisner-section-sr}. A
simplicial complex whose geometric
realization is homeomorphic to a sphere is
called a triangulated sphere. The question of
the next sections is the converse of the
previous paragraph: which triangulated
spheres are boundary complexes of convex
polytopes?

\section{Polytopality and nonrealizable
abstract symmetries}
\label{section-polytopality-symmetries}

\noindent
The next four sections follow
\cite{HollebenYang2026}. They explain what
it means for a triangulated sphere to be
the boundary of a convex polytope, and why
a symmetry of the triangulation need not
be a symmetry of any convex realization.

\section{What polytopality asks}
\label{subsection-polytopality}

\noindent
Source: Thiago Holleben and Yirong Yang,
\emph{Polytopal Bier spheres and
nonrealizable central symmetries}
\cite{HollebenYang2026} (2026 preprint).

Start with an abstract triangulated sphere:
we know which vertices form faces, but
have not assigned coordinates. Can we
place its vertices so that these are
exactly the faces of a convex solid?
A simplicial $d$-sphere is
\emph{polytopal} if it is isomorphic to
the boundary of a convex simplicial
$(d+1)$-polytope. Convexity concerns the
solid, not its boundary: every segment
between points of the solid stays inside.
A nonconvex drawing does not prove
nonpolytopality; every possible convex
realization must be ruled out.

Nonpolytopality starts in sphere dimension
$d=3$, hence ambient dimension four.
A triangulated circle is a polygon.
For a triangulated $2$-sphere, its graph
is planar and $3$-connected; Steinitz's
theorem gives a convex $3$-polytope with
that graph and the same faces. See
\href{https://arxiv.org/abs/2608.07233}
{Holleben--Yang, Introduction}.
In higher dimensions, being a topological
sphere no longer guarantees compatible
convex coordinates: all prescribed
facets must simultaneously lie in
supporting hyperplanes, with all other
vertices strictly on the inner side.
These geometric constraints can conflict.

\section{Realizing abstract symmetries}
\label{subsection-realizing-symmetries}

\noindent
What's the relationship between
realizability of symmetries and
polytopality?

\noindent
A vertex permutation $\varphi$ preserving
faces is an abstract symmetry. In a
realization with vertices $p_v$, it is
realized affinely when an affine
bijection $A$ satisfies
$$
A(p_v)=p_{\varphi(v)}.
$$
A symmetry missing from one realization
may appear in another. A nonrealizable
symmetry appears in none of the convex
realizations. This still does not imply
that the sphere itself is nonpolytopal.

Central symmetry geometrically sends
$x$ to $2c-x$ about a centre $c$.
Combinatorially it pairs vertices by a
face-preserving involution fixing no
nonempty face setwise.

For the Stanley--Reisner scheme,
every abstract symmetry acts by permuting
variables, because it preserves $I_K$.
Convex realizability is not needed for
this algebraic action.

\section{Bier spheres and an infinite
family}
\label{subsection-bier-spheres}

\noindent
Bier's construction turns a complex into
a sphere by using two copies of its
vertices. Choose a face from the first
copy and a face from a complementary
complex in the second. You may combine
them provided no vertex is chosen in both
copies. Remarkably, the resulting complex
is a sphere even when the input is not.

\begin{definition}
\label{definition-bier-sphere}
Let $V=\{1,\ldots,n\}$ and let $K$ be a
proper subcomplex of $2^V$ containing
$\varnothing$. Its \emph{Alexander dual}
is
$$
K^\vee=\{B\subseteq V:V\setminus B\notin K\}.
$$
On two copies $V^+$ and $V^-$, the
\emph{Bier sphere} has faces
$$
\begin{gathered}
A^+\cup B^-,\\
A\in K,\quad B\in K^\vee,\quad
A\cap B=\varnothing.
\end{gathered}
$$
\end{definition}

\noindent
Reference:
\href{https://arxiv.org/abs/2608.07233}
{Holleben--Yang, Section 2.2}.
Bier's theorem says this is an
$(n-2)$-sphere with at most $2n$ vertices;
some singleton sets may be absent.

With this construction defined, we can
state the failure of symmetry realization
for spheres which are nevertheless
polytopal.

\begin{theorem}[Holleben--Yang]
\label{theorem-polytopal-nonrealizable}
There exist polytopal Bier spheres whose
combinatorial central symmetry cannot be
realized affinely in any convex
realization.
\end{theorem}

\noindent
Reference:
\href{https://arxiv.org/abs/2608.07233}
{Holleben--Yang, Theorem 1.2 and
Lemma 3.10}.
This disproves the tempting implication
from nonrealizability of a symmetry to
nonpolytopality of the sphere.
For their twelve-vertex examples,
Proposition 3.4 rules out central
symmetry geometrically, while Theorem 1.4
establishes polytopality.

The following result supplies an infinite
family when the dimension is allowed to
grow.

\begin{theorem}[Holleben--Yang]
\label{theorem-bier-infinite-family}
For every even $n\geq6$, there are
combinatorially centrally symmetric Bier
$(n-2)$-spheres with $2n$ vertices which
admit no centrally symmetric convex
realization.
\end{theorem}

\noindent
This is a consequence of
\href{https://arxiv.org/abs/2608.07233}
{Holleben--Yang, Theorem 1.1}.
The vertex counts are $12,16,20,\ldots$.
The theorem does not assert that these
spheres are nonpolytopal. All Bier spheres
with at most twelve vertices are
polytopal by Theorem 1.4 of that preprint.

\section{Where the BEK example fits}
\label{subsection-bek-example}

\noindent
BEK stands for
Bokowski--Ewald--Kleinschmidt.
Their example is a polytopal $3$-sphere
with ten vertices, and therefore the
boundary of a convex $4$-polytope.
Its combinatorial central symmetry is
$$
(1\,10)(2\,5)(3\,6)(4\,9)(7\,8).
$$
No convex realization realizes this
permutation affinely. See
\href{https://arxiv.org/abs/2608.07233}
{Holleben--Yang, Section 3.3, Proposition 3.8
and Lemma 3.10}.
It is the earlier counterexample to the
implication discussed above; the Bier
construction supplies further examples.

There is also a topological consequence:
for a polytopal sphere with nonrealizable
central symmetry, the space of convex
realizations modulo affine changes of
coordinates is noncontractible;
see \href{https://arxiv.org/abs/2608.07233}
{Holleben--Yang, Proposition 3.14}.

\section{The Grünbaum-Sreedharan sphere}
\label{section-grunbaum-sreedharan-sphere}

\noindent
Grünbaum and Sreedharan enumerated the
simplicial convex polytopes of dimension $4$
with $8$ vertices \cite{gru}. In the course
of the enumeration they found a triangulation
$\mathcal{M}$ of the sphere of dimension $3$,
with $8$ vertices and $20$ tetrahedra, that
is not the boundary complex of any convex
polytope. So non-polytopal spheres exist
already with $8$ vertices in dimension $3$.

The Stanley-Reisner ideal of $\mathcal{M}$ is
written down in Stanley-Reisner schemes,
Section
\ref{stanley-reisner-section-grunbaum-sphere}.

\section{6-vertex triangulation of RP2}
\label{section-6-vertex-triangulation-of-rp2}

\noindent
Here I read \cite{KuhnelBanchoff1983}.

\medskip\noindent
The icosahedron quotiented by antipodal
map gives a triangulation of 6 vertices
of $\mathbb{R}P^2$.
This can be embedded in $\mathbb{R}^5$ 
inside the 5-simplex $\Delta^5$,
which has 6 vertices.
This embedding has the properties
of symmetry (icosahedral$/-\text{Id}$),
duality (three vertices of $\Delta^5$
determine
a triangle of $\mathbb{R}P^2_6$ if and only
if
the remaining do not),
tightness (too complicated) and
secant tangency (a chord joining
two points of $\mathbb{R}P^2$ is contained
in a boundary 4-simplex of $\Delta^5$.

Then consider Veronesse embedding
$[x_0:x_1:x_2]\mapsto
[x_0^2:x_1^2:x_2^2:x_0x_1:x_0x_2:x_1x_2]$.
Like $\mathbb{R}P^2_6$ is contained in the
boundary of $\Delta^5$,
this embedding is contained in $S^4$.
It also satisfies properties of symmetry (of
$\mathbb{R}^3$),
duality, tightness and secant-tangency.

Each of the third and fourth properties
characterizes $\mathbb{R}P^2_v$ among
the smooth embeddings of $\mathbb{R}P^2$
in $\mathbb{R}^5$ not lying in an affine
hyperplane.
``Pohl and Kuipar have in fact shown that
any {\it topological} embedding of the real
projective plane into $\mathbb{R}^5$ not
lying in
a hyperplane which has the tightness property
alone must either be $\mathbb{R}P^2_v$ or
$\mathbb{R}P^2_v$.''

With respect to the complex projective plane,
the situation has been less satisfactory.

\bibliography{my}
\bibliographystyle{amsalpha}

\end{document}
